\section{Группы Кремоны и локально нильпотентные дифференцирования}

Ниже $k$ --- поле характеристики нуль, $A=k[x,y,z]$. Все утверждения в
этом разделе доказаны элементарно; мотивы, $K$-теория и гомологическая
алгебра не используются. Общие приёмы не объявляются новыми. В частности,
критерий подъёма параметра и устранение полюсов рациональных сдвигов
сопоставляются с работой Drew Lewis,
\emph{A unified approach to embeddings of a line in 3-space},
\url{https://arxiv.org/abs/2203.08178}: Lemma~1, Theorem~2 и Corollary~4.

\subsection{Точный критерий внутри класса переносов}

Пусть $(x,H,W)$ --- автоморфизм $\mathbb A^3_k$, сохраняющий $x$, и
\[
 d=H_yW_z-H_zW_y\in k^*,\qquad
 \beta=(x,p(x,y)+m(x)z,y),\qquad 0\ne m\in k[x].
\]
Обозначим
\[
 h=H\circ\beta,\quad w=W\circ\beta,\quad
 B=k[x]/(m),\quad G(y)=\overline{H(x,p(x,y),y)}\in B[y].
\]
Черта означает редукцию по $m$. В случае $m\in k^*$ перенос уже является
полиномиальным автоморфизмом; содержательный случай --- $\deg m>0$.

\paragraph{Утверждение.}
Дифференцирование
\[
 \delta_h=h_z\partial_y-h_y\partial_z
\]
всегда локально нильпотентно. Следующие четыре условия эквивалентны:
\begin{enumerate}
\item $G$ является координатой однопеременного кольца $B[y]$, то есть
подстановка $y\mapsto G(y)$ есть $B$-автоморфизм;
\item существуют $q\in A$ такие, что $(x,h,q)$ --- полиномиальный
автоморфизм;
\item у $\delta_h$ существует слайс: $\delta_h(s)=1$ для некоторого $s\in A$;
\item $\delta_h$ не имеет геометрических неподвижных точек, эквивалентно
$(h_y,h_z)=A$.
\end{enumerate}
В условии~1, если $G=\sum g_jy^j$, требуется и достаточно, чтобы $g_1$
был обратим в $B$, а все $g_j$ при $j\ge2$ были нильпотентны. Коэффициент
$g_0$ произволен. Поэтому проверка всех корней $m$ должна устанавливать
\emph{аффинность всей зависимости от $y$}, а не только ненулевую производную
в одной выбранной точке.

\paragraph{Доказательство локальной нильпотентности.}
В координатах $(x,H,W)$ дифференцирование
$\delta_H=H_z\partial_y-H_y\partial_z$ равно $-d\partial_W$.
После обращения $m$ подстановка $\beta^*$ является автоморфизмом и
\[
 \delta_h=-m\,\beta^*\delta_H(\beta^*)^{-1}.
\]
Множитель $m$ принадлежит ядру, поэтому правая часть локально
нильпотентна на $A_m$. Левая часть сохраняет $A$, а вложение $A\hookrightarrow
A_m$ инъективно; следовательно, её ограничение на $A$ также локально
нильпотентно. Это верно \emph{даже при невыполнении условия центра}.

\paragraph{Доказательство $1\Rightarrow2,3$.}
Пусть $f(T)\in B[T]$ --- обратная подстановка к $G$. Положим
\[
 J(y)=\overline{W(x,p(x,y),y)},\qquad \bar\varphi(T)=J(f(T)),
\]
и выберем произвольный полиномиальный подъём $\varphi\in k[x,T]$.
Тогда
\[
 q=\frac{w-\varphi(x,h)}{m}\in A.
\]
Обратное отображение строится без предположения об отсутствии кратных
корней. Пусть $(x,Y(x,T,V),Z(x,T,V))$ --- обратное к $(x,H,W)$.
При $V=mQ+\varphi(x,T)$ положим
\[
 y=Z(x,T,V),\qquad z=\frac{Y(x,T,V)-p(x,Z(x,T,V))}{m}.
\]
В кольце $B[T]$ имеем
\[
 (Y(x,T,\bar\varphi(T)),Z(x,T,\bar\varphi(T)))
   =(\bar p(x,f(T)),f(T)).
\]
Значит, последнее деление также точное. Обе композиции равны тождеству
после обращения $m$, следовательно, и в полиномиальном кольце. Кроме того,
\[
 \delta_h(w)=md,\qquad \delta_h(q)=d,
 \qquad \ker\delta_h=k[x,h].
\]
Таким образом, $q/d$ --- слайс. Именно необходимость проверки обратного
деления защищает от ошибочного перехода от бирациональности к автоморфизму.

\paragraph{Остальные импликации.}
Если $(x,h,q)$ --- автоморфизм, то его якобиан есть ненулевая константа,
так что $\delta_h(q)$ --- ненулевая константа: $2\Rightarrow3$.
Слайс несовместим с неподвижной точкой, откуда $3\Rightarrow4$.
Вне $m=0$ поле $\delta_h$ нигде не обращается в нуль, поскольку сопряжено
с ненулевой кратной постоянного поля. На геометрическом слое $x=\alpha$,
где $m(\alpha)=0$, оно равно
\[
 \delta_h\big|_{x=\alpha}=-G_\alpha'(y)\partial_z.
\]
Над алгебраическим замыканием полином $G_\alpha'$ не имеет корней тогда
и только тогда, когда он постоянен и отличен от нуля. В характеристике
нуль это равносильно тому, что $G_\alpha$ аффинен с ненулевым наклоном.
Для всех корней $m$ это ровно обратимость $g_1$ и нильпотентность высших
коэффициентов в $B$. Наконец, такая подстановка обратима: после обращения
линейной части остаток имеет коэффициенты в нильпотентном идеале, и
последовательное подставление заканчивается через конечное число шагов.
Получаем $4\Rightarrow1$.

\paragraph{Граница вывода.}
Это необходимый и достаточный критерий \emph{относительной} координаты
$(x,h)$ в указанном классе переносов. Он не является необходимым критерием
выпрямления произвольной кривой и здесь не утверждается как критерий того,
что $h$ есть абсолютная координата без сохранения $x$.
Если условие центра не выполнено, отвергаются именно этот перенос и слайс
именно этого $\delta_h$. Другие координаты и другие LND остаются возможны.
Без условия центра нельзя автоматически писать $\ker\delta_h=k[x,h]$:
например, $H=Y$, $p=xy$, $m=x$ дают $h=x(y+z)$ и
$\ker\delta_h=k[x,y+z]\supsetneq k[x,h]$.

\subsection{Конструктивная лемма о треугольных LND}

Рассмотрим $k[x,v,y]$ и
\[
 D=P(x)\partial_v+Q(x,v)\partial_y,\qquad 0\ne P\in k[x].
\]
Выберем $B_0(x,v)$ с $(B_0)_v=Q$ и положим
\[
 H=B_0(x,v)-P(x)y.
\]
Предположим, что в $k[x]/(P)$ постоянный по $v$ коэффициент $Q$ обратим,
а все остальные коэффициенты $Q$ нильпотентны. Тогда
$\bar B_0$ является автоморфизмом $(k[x]/(P))[v]$. Если $f(x,T)$ ---
полиномиальный подъём его обратного, то
\[
 W=\frac{v-f(x,H)}{P(x)}\in k[x,v,y],\qquad DH=0,\quad DW=1.
\]
Полиномиальность $W$ следует из $H\equiv B_0(v)\pmod P$. Обратные
формулы имеют вид
\[
 v=f(x,H)+P(x)W,\qquad
 y=\frac{B_0(x,f(x,H)+P(x)W)-H}{P(x)}.
\]
Второй числитель также делится на $P$, поскольку
$B_0(x,f(x,H))\equiv H\pmod P$. Поэтому $(x,H,W)$ --- автоморфизм,
\[
 k[x,v,y]=k[x,H,W],\qquad D=\partial_W,
 \qquad \ker D=k[x,H].
\]
Над алгебраически замкнутым полем гипотеза леммы эквивалентна отсутствию
неподвижных точек у этого треугольного LND: над каждым корнем $P$
полином $Q(\alpha,v)$ должен быть ненулевой константой.

В частном случае $Q=Q_1(x)v+R(x)$ критерий принимает простой вид
\[
 Q_1\in\sqrt{(P)},\qquad \gcd(P,R)=1.
\]
Для базового дифференцирования репозитория это объясняет роль
$P=(1-a)x^5$, $Q_1=2(1-a)x$, $R=-1+ax^2-ax^4$: старшая поправка
нильпотентна по модулю $x^5$, а линейная часть обратима. Это конструктивная
форма классического механизма, а не заявка на новую общую классификацию LND.

\paragraph{Небольшой полностью проверенный пример с кратным корнем.}
Для $m=x^2$, $p=y+xy^2$, $H=Y$, $W=Z$ центр имеет обратный
$f(T)=T-xT^2\pmod{x^2}$. Получаем
\[
 h=y+xy^2+x^2z,\qquad
 q=x^3z^2+2x^2y^2z+xy^4+2xyz+2y^3-z.
\]
Тогда $\delta_h=x^2\partial_y-(1+2xy)\partial_z$,
$\delta_h(h)=0$, $\delta_h(q)=1$, и обратные формулы к $(x,h,q)$ суть
\[
 y=T-xT^2+x^2Q,
\]
\[
 z=-Q^2x^3+2QT^2x^2-2QTx-Q-T^4x+2T^3.
\]
Обе композиции, равенства для LND и якобиан $-1$ проверены точно
скриптом \texttt{research/cremona\_lnd/verify\_example.py}.
Напротив, для $h=y^2+xz$ существует LND
$x\partial_y-2y\partial_z$, но она имеет неподвижную прямую $x=y=0$
и потому не имеет слайса. Это простой пример того, почему одного
построения LND недостаточно.

\subsection{Что даёт группа Кремоны и где возникают полюса}

Пусть $\gamma(t)=(f(t),g(t),h_0(t))$ --- замкнутое вложение. После общей
линейной замены двух координат можно добиться
$k(f,g)=k(t)$: это обычная теорема о примитивном элементе для конечного
сепарабельного расширения $k(t)/k(f)$. Поэтому найдётся
$r(x,y)\in k(x,y)$ с $r(f(t),g(t))=t$.
Положим
\[
 w=z-h_0(r(x,y))+r(x,y),\qquad
 F(x,y,z)=(x-f(w),y-g(w),w).
\]
Это бирациональное преобразование с явным обратным
\[
 x=X+f(s),\quad y=Y+g(s),\quad
 z=s-r(x,y)+h_0(r(x,y)).
\]
На плотном открытом подмножестве кривой $F\gamma(t)=(0,0,t)$.
Таким образом, бирациональное выпрямление здесь имеет элементарную
конструкцию. Оно ещё не даёт отображения, регулярного во всех точках
кривой, тем более глобального полиномиального автоморфизма.

Если числители и знаменатели построенного преобразования лежат в
$A_m$, задача полиномиальности сводится к точной делимости:
\[
 \frac{N}{m^e}\in A
 \quad\Longleftrightarrow\quad
 \operatorname{ord}_\pi(N)\ge e\operatorname{ord}_\pi(m)
 \quad\hbox{для всех неприводимых }\pi\mid m.
\]
При $m=\prod_\alpha(x-\alpha)^{r_\alpha}$ это равносильно занулению
всех $x$-производных числителя порядков $0,\ldots,er_\alpha-1$ при
$x=\alpha$. Проверка только значений в корнях недостаточна для кратных
корней. Требования применяются к прямому и обратному отображениям.

Предупреждение о LND существенно: рациональное сопряжение постоянного
поля может стать полиномиальным полем, не будучи локально нильпотентным.
Например, $\partial_u$ при $u=-1/x$ превращается в $x^2\partial_x$.
Безопасный вариант таков: если сопряжение является автоморфизмом $A_m$
над $k[x,m^{-1}]$, то сопряжённые $\partial_y,\partial_z$ локально
нильпотентны на $A_m$ и аннулируют $m$. Умножение на достаточно большую
степень $m$ делает их полиномиальными LND на $A$; два поля остаются
коммутирующими. Но возникающее вырождение на $m=0$ может уничтожить
слайсы. Поэтому очистка знаменателей LND сама по себе не решает задачу
выпрямления.

\subsection{Подъём параметра, стабилизация и действия на дополнении}

Для идеала $I=I(\gamma)$ замкнутость вложения означает наличие
$R\in A$ с $R\circ\gamma=t$. Выпрямляемость равносильна тому, что класс
$R+I$ содержит \emph{координату} $S$ кольца $A$. Действительно, если
$(S,U,V)$ --- координаты и $S\gamma=t$, то
\[
 (S,U-U(\gamma(S)),V-V(\gamma(S)))
\]
выпрямляет кривую; здесь $U(\gamma(t)),V(\gamma(t))$ сначала рассматриваются
как однопеременные полиномы, после чего подставляется $S$.
Обратное следует из первой компоненты выпрямляющего автоморфизма.
Это известный критерий; см. Lewis, Lemma~1.

После добавления \emph{одной} переменной любое такое вложение
выпрямляется явно. На $\mathbb A^4$ последовательно применяются
\[
 (x,y,z,u)\longmapsto(x,y,z,u+R(x,y,z))
\]
и
\[
 (x,y,z,u)\longmapsto(x-f(u),y-g(u),z-h_0(u),u).
\]
Они переводят $(\gamma(t),0)$ в $(0,0,0,t)$ и являются произведениями
элементарных полиномиальных автоморфизмов. Этот факт не позволяет убрать
добавленную переменную; именно такой переход и требует отдельного
доказательства в размерности три.

Если выпрямитель уже известен, обозначим его координаты $(S,U,V)$ так,
что $I=(U,V)$ и $S\gamma=t$. Сопряжение трёх частных производных даёт
коммутирующие LND $D_S,D_U,D_V$ с
\[
 D_i(S,U,V)=e_i.
\]
В частности, $D_S$ имеет слайс $S$, ядро $k[U,V]$ и орбиту $\gamma$.
Явное действие записывается через выпрямитель $\Phi$:
\[
 \lambda\cdot p=\Phi^{-1}(S(p)+\lambda,U(p),V(p)).
\]
Это даёт полноценное свободное алгебраическое действие $\mathbb G_a$,
а не только касательное поле вдоль кривой. Обратное утверждение о
\emph{произвольной} орбите произвольного LND здесь не используется.

Обычный инвариант Макар--Лиманова кольца регулярных функций дополнения
не различает рассматриваемые вложения. Для любой кривой
$C\subset\mathbb A^3$ коразмерности два
\[
 \Gamma(\mathbb A^3\setminus C,\mathcal O)=k[x,y,z].
\]
Элементарное объяснение: неприводимый множитель знаменателя рациональной
функции задаёт дивизор; такой дивизор не может целиком содержаться в $C$,
поэтому регулярность на дополнении исключает каждый такой множитель.
Следовательно, ML этого кольца всегда равен $k$.

Для сохранения информации о вложении можно рассматривать именно пару
$(A,I)$ и пересечение
\[
 \mathrm{ML}(A,I)=
 \bigcap_{\substack{D\in\operatorname{LND}(A)\\D(I)\subset I}}\ker D.
\]
Это определение фиксирует, что соответствующие действия сохраняют
удаляемую кривую. Для выпрямляемой кривой оно также равно $k$:
достаточно трёх сохраняющих $I$ LND
$D_S$, $U D_V$, $V D_U$, чьи ядра имеют пересечение $k$.
Таким образом, нетривиальность этого относительного пересечения была
бы препятствием, но его тривиальность не объявляется достаточной для
выпрямления. Вычисление всех сохраняющих $I$ LND --- отдельная трудная
задача; простая проверка конечного списка полей не устанавливает
нетривиальность полного пересечения.

\paragraph{Дополнительные первоисточники.}
K. Masuda, \emph{Homogeneous locally nilpotent derivations having slices
and embeddings of affine spaces}, J. Algebra 321 (2009), 1719--1733,
\url{https://doi.org/10.1016/j.jalgebra.2008.12.012}, устанавливает
связь вложений с определёнными однородными коммутирующими LND со
слайсами; важно сохранять условия на градуировку.
S. Kaliman, \emph{Free $\mathbb C_+$-actions on $\mathbb C^3$ are
translations}, \url{https://arxiv.org/abs/math/0207156}, даёт общий
результат о свободных действиях в размерности три. Его глубокая теорема
не нужна для доказанных выше специальных критериев.
